% R008 — Guide to Cultivating Complex Analysis (Bahasa Indonesia)
% Reader-facing translation unit: source/ca-v1.9/ca.tex lines 1384–1398.
% Source release: v1.9 (commit a4d25d9ba37a486a5a020219eb8bd4e6602fd14c).
% Translation provenance: OpenAI Codex gpt-5.6-sol, Ultra, acting on the user's request.
% Preserve source equations, notation, and labels.

\subsection{Sifat pemetaan eksponensial}

Mari kita lihat bagaimana eksponensial memetakan bidang kompleks.
Identitas $e^{z+w} = e^z e^w$ menyiratkan bahwa eksponensial
tidak pernah bernilai nol (latihan).
Dari sifat-sifat koordinat polar yang telah diketahui
dan eksponensial real, diperoleh bahwa eksponensial kompleks bersifat
surjektif ke $\C \setminus \{ 0 \}$.
Eksponensial kompleks tidak satu-ke-satu---setiap nilai memiliki
tak hingga banyak prapeta. Untuk sembarang bilangan bulat $k$,
\begin{equation} \label{eq:exponentialidentities}
e^{z+ik2\pi} =
e^{z} e^{ik2\pi} =
e^z .
\end{equation}
